\documentclass[11pt]{article}
\usepackage[a4paper,margin=26mm]{geometry}
\usepackage{amsmath,amssymb,amsthm,booktabs,array,longtable}
\usepackage[T1]{fontenc}
\usepackage{lmodern}
\usepackage[colorlinks=true,urlcolor=blue,citecolor=blue,linkcolor=blue]{hyperref}
\hypersetup{pdftitle={Prime-shift extensions of Kaplansky-Riordan's menage identities: Mathar's recurrence and van Hoeij's congruences},pdfauthor={Haobo Ma and Wenlin Zhang}}
\newtheorem{theorem}{Theorem}
\newtheorem{lemma}{Lemma}
\newtheorem{corollary}{Corollary}
\newcommand{\coeff}[2]{[#1^{#2}]}
\title{Prime-shift extensions of Kaplansky--Riordan's m\'enage identities:\protect\\
Mathar's recurrence and van Hoeij's congruences}
\author{Haobo Ma\thanks{ChronoAI Pte Ltd, Singapore.}\and
Wenlin Zhang\thanks{National University of Singapore; ChronoAI Pte Ltd, Singapore.}}
\date{15 November 2026}
\begin{document}
\maketitle
\begin{abstract}
Extending Kaplansky and Riordan's m\'enage identities (1946,
pp.~121--122, equations (25) and (33)), we use Jovovi\'c's factorial-series
generating function to prove the conjectures of Mathar (recurrence) and
van Hoeij (prime-shift congruence) in OEIS A000159, revision 54, first and third
FORMULA lines, dated 2 November 2015 and 10 June 2019.
For the hit numbers $u_h(n)=[t^h]U_n(t)$, the stronger uniform result is
$u_h(n)+2u_h(n+p)+u_h(n+2p)\equiv0\pmod p$ for every prime $p$,
$1\leq h<p$ and $n\geq0$; the zero-hit case holds for $n\geq1$.
It settles the neighboring two-, four- and five-hit congruences.
We also prove prime-square antiperiodicity for three hits, give exact
recurrence certificates for the neighboring diagonals, and correct the
one-hit indexing formula by identifying it with a classical identity.
\end{abstract}

\section{The sequence and the statements}
The main observation is that factorial truncation modulo $p$ gives the
same denominator for every hit degree below $p$. It proves several
prime-shift statements at once, while the characteristic-zero generating
function proves the recurrences. Kaplansky--Riordan's classical
differentiation identities supply the one-hit correction and the
two-hit/three-hit relation; those identities are reproved with their
original attribution, not presented as new.

Mathar's recurrence concerns the diagonal
$a(n)=\mathrm{A058087}(n,n-3)$ of the triangle of m\'enage hit
polynomials \cite{a159,a58087}. We prove that recurrence and show
that the same generating function proves van Hoeij's prime-shift
congruence. Throughout, generating functions are formal power series;
no analytic convergence is required.

Write $U_0(t)=1$ and, for $n\geq1$, put
\begin{equation}\label{eq:rook}
 U_n(t)=\sum_{j=0}^{n}\frac{2n}{2n-j}
           \binom{2n-j}{j}(n-j)!(t-1)^j.
\end{equation}
This is the m\'enage hit polynomial of Kaplansky and Riordan
\cite[p.~121, equation (25)]{kr}; it is the formula recorded in
OEIS A058087 \cite[FORMULA, Orrick, 9 August 2020]{a58087}.
That triangle lists coefficients in descending degree, so
$a(n)=\coeff{t}{3}U_n(t)$ for $n\geq3$.
Set $a(0)=a(1)=a(2)=0$ and
$A(x)=\sum_{n\geq0}a(n)x^n$.
Put $u_h(n)=[t^h]U_n(t)$ and $C_h(x)=\sum_{n\geq0}u_h(n)x^n$,
so $a=u_3$ and $A=C_3$. The first six terms in the source domain are
\[
 (a(3),\ldots,a(8))=(2,8,20,152,994,7888).
\]
\begin{theorem}[Uniform prime shifts]\label{thm:uniform}
For every prime $p$, every integer $1\leq h<p$, and every $n\geq0$,
\[
 u_h(n)+2u_h(n+p)+u_h(n+2p)\equiv0\pmod p.
\]
The same congruence holds for $h=0$ and $n\geq1$.
\end{theorem}

For $0\leq j\leq6$, define $c_j(n)$ by
\begin{center}
\begin{tabular}{cl}
\toprule
$j$ & $c_j(n)$\\
\midrule
0 & $2(-252307n+1041077)$\\
1 & $504614n^2-3362985n+5118150$\\
2 & $1280831n^2-7397886n+6461565$\\
3 & $746598n^2-2913543n-1336090$\\
4 & $-405481n^2+6175011n-15469320$\\
5 & $-375862n^2+4098537n-8846430$\\
6 & $2(-187931n+560630)$\\
\bottomrule
\end{tabular}
\end{center}

\begin{theorem}\label{thm:rec}
For every $n\geq9$,
\begin{equation}\label{eq:rec}
 \sum_{j=0}^{6}c_j(n)a(n-j)=0.
\end{equation}
With the zero extension above, the same identity holds for $n=6,7,8$.
\end{theorem}

\begin{theorem}\label{thm:cong}
For every prime $p\ne3$ and every $n\geq0$,
\begin{equation}\label{eq:cong}
 a(n)+2a(n+p)+a(n+2p)\equiv0\pmod p.
\end{equation}
\end{theorem}

\section{A factorial-series expression}
Let $F(z)=\sum_{k\geq0}k!z^k$, let
$Z=x/(1+x)^2$, and let $B(x)=F(Z)$.
Composition is well defined because $Z$ has zero constant term.
Define the integer polynomials
\begin{align*}
 r(x)&=-4x^5+6x^4+10x^3+8x^2+6x+4,\\
 s(x)&=4x^5-2x^4-16x^3-14x^2-10x-4,\\
 d(x)&=3(1+x)^2.
\end{align*}

\begin{lemma}\label{lem:gf}
The series $A$ satisfies
\begin{equation}\label{eq:affine}
 d(x)A(x)=r(x)B(x)+s(x).
\end{equation}
\end{lemma}
\begin{proof}
For $n\geq1$ and $0\leq j\leq n$,
\[
 \frac{2n}{2n-j}\binom{2n-j}{j}
 =\binom{2n-j}{j}+\binom{2n-j-1}{j-1},
\]
where the second term is zero when $j=0$.
The identity follows from
$j\binom{2n-j}{j}=(2n-j)\binom{2n-j-1}{j-1}$.
In particular, every coefficient of $U_n$ is an integer.
Substitute $n=k+j$ into \eqref{eq:rook}, sum over $n$, and put
$u=x(t-1)$. The binomial-series identities give
\begin{align}
 \sum_{n\geq0}U_n(t)x^n
 &=\sum_{k\geq0}k!x^k
   \sum_{j\geq0}\left\{\binom{2k+j}{j}
           +\binom{2k+j-1}{j-1}\right\}u^j\notag\\
 &=\sum_{k\geq0}k!x^k\frac{1+u}{(1-u)^{2k+1}}.
 \label{eq:bivariate}
\end{align}
For $k=j=0$, the brace is interpreted as $1$, consistently with $U_0=1$.
All rearrangements are valid coefficientwise: only finitely many terms
contribute to a fixed power of $x$.
Equation~\eqref{eq:bivariate} is the generating function recorded by
Vladeta Jovovi\'c in A058087 (first FORMULA, 14 December 2009),
after the substitution $(x,y)\mapsto(xt,t^{-1})$ reverses the row order
\cite{a58087}. We use this classical input to derive the scalar
recurrence and the uniform congruences.

Write $\alpha_k=\binom{2k+3}{3}$ and $\beta_k=\binom{2k+2}{2}$.
The coefficient of $t^3$ in the $k$th summand in
\eqref{eq:bivariate} is
\[
 \frac{k!x^{k+3}}{(1+x)^{2k+4}}
       \bigl((1-x)\alpha_k+(1+x)\beta_k\bigr).
\]
Hence, with $\theta=z\,d/dz$,
\begin{equation}\label{eq:theta}
 3A=\frac{x^3}{(1+x)^4}
 \left[(4-4x)\theta^3F+(18-6x)\theta^2F
                  +(20-2x)\theta F+6F\right]_{z=Z}.
\end{equation}
Indeed,
$3(\alpha_k+\beta_k)=4k^3+18k^2+20k+6$ and
$3(-\alpha_k+\beta_k)=-4k^3-6k^2-2k$.
Direct coefficient comparison in $F$ yields
\[
 z^2F'(z)+(z-1)F(z)=-1.
\]
Differentiation, or iteration of $\theta$, now gives
\begin{align*}
 z\theta F&=(1-z)F-1,\\
 z^2\theta^2F&=(z^2-3z+1)F+2z-1,\\
 z^3\theta^3F&=(-z^3+7z^2-6z+1)F-3z^2+5z-1.
\end{align*}
Insert these three identities into \eqref{eq:theta} and use
$Z=x/(1+x)^2$. Collecting the coefficient of $F(Z)$ and the
constant term gives, respectively, $r/[3(1+x)^2]$ and
$s/[3(1+x)^2]$. This proves \eqref{eq:affine}.
The calculation can be made in $\mathbb Q((x))$; its final identity
lies in $\mathbb Q[[x]]$.
\end{proof}

\section{The differential identity and the recurrence}
The chain rule applied to the factorial-series equation gives
\begin{equation}\label{eq:B}
 x^2(1+x)B'=(1-x^3)B+(x^2-1)(1+x).
\end{equation}
For clarity, define rational functions
\[
 R=\frac r d,\qquad S=\frac s d,\qquad
 T=\frac{1-x^3}{x^2(1+x)},\qquad V=\frac{x^2-1}{x^2}.
\]
Thus $A=RB+S$ and $B'=TB+V$ in the Laurent-series field.
Set
\[
 R_1=R'+RT,\quad S_1=S'+RV,\quad
 R_2=R_1'+R_1T,\quad S_2=S_1'+R_1V.
\]
Then $A'=R_1B+S_1$ and $A''=R_2B+S_2$.
The following polynomials give a differential equation for $A$:
\begin{align*}
 P_2={}&-375862x^7-405481x^6+746598x^5
                         +1280831x^4+504614x^3,\\
 P_1={}&-375862x^7-35945x^6+2525682x^5+2312643x^4\\
       &\hspace{16mm}-993731x^3-1849143x^2-504614x,\\
 P_0={}&-1133912x^6+2249705x^5+2743028x^4-3357337x^3\\
       &\hspace{16mm}-3210883x^2+2259779x+2082154,\\
 H={}&1136624x^3-10348x^4+1516100x^5.
\end{align*}
Explicit rational identities are
\begin{equation}\label{eq:cert}
 \begin{split}
 P_2R_2+P_1R_1+P_0R&=0,\\
 P_2S_2+P_1S_1+P_0S&=H.
 \end{split}
\end{equation}
These identities can be verified by substitution of the displayed
$r,s,d,T,V$ and polynomial expansion after clearing denominators.
They imply
\begin{equation}\label{eq:ode}
 P_2A''+P_1A'+P_0A=H.
\end{equation}

\begin{proof}[Proof of Theorem~\ref{thm:rec}]
Let $D=x\,d/dx$. Substitution of the seven displayed $c_j$ gives the
operator identity
\[
 \sum_{j=0}^6x^jc_j(D+j)
       =P_2\frac{d^2}{dx^2}+P_1\frac{d}{dx}+P_0.
\]
To check the equality, write $c_j(v)=u_jv^2+v_jv+w_j$.
The three coefficients on the right are obtained on the left as
\[
 x^2\sum_j u_jx^j,\qquad
 x\sum_j\bigl((2j+1)u_j+v_j\bigr)x^j,\qquad
 \sum_j(u_jj^2+v_jj+w_j)x^j,
\]
using $D^2=x^2(d/dx)^2+x(d/dx)$.
For $n\geq6$, the coefficient of $x^n$ on the left, applied to $A$,
is $\sum_j c_j(n)a(n-j)$: $D+j$ acts on $x^{n-j}$ with eigenvalue $n$.
The coefficient of $x^n$ in $H$ vanishes because $\deg H=5$.
Equation~\eqref{eq:ode} therefore proves the claim for all $n\geq6$,
including Mathar's full source domain $n\geq9$.
\end{proof}

Since $c_0(n)\ne0$ for every integer $n\geq9$, the recurrence and the
six values $a(3),\ldots,a(8)$ determine the sequence on its source domain.

\section{The prime-shift congruence}
\begin{proof}[Proof of Theorem~\ref{thm:uniform}]
Work in $\mathbb F_p[t][[x]]$ and set $u=x(t-1)$.
Factorial truncation in \eqref{eq:bivariate} gives
\begin{equation}\label{eq:uniformnumerator}
 (1-u)^{2p}\sum_{n\geq0}U_n(t)x^n
 =N_p(x,t):=\sum_{k=0}^{p-1}k!x^k(1+u)(1-u)^{2p-2k-1}.
\end{equation}
Each summand has $x$-degree at most $2p-k$, hence $\deg_xN_p\leq2p$.
By Frobenius,
\[
 (1-u)^{2p}=(1+x^p-x^pt^p)^2.
\]
For $0\leq h<p$, the coefficient of $t^h$ on the left of
\eqref{eq:uniformnumerator} is therefore $(1+x^p)^2C_h(x)$.
Only $k=0$ contributes to $[x^{2p}]N_p$, and that coefficient is
$-(t-1)^{2p}=-(t^{2p}-2t^p+1)$.
Its $t^h$ coefficient is zero for $1\leq h<p$.
Thus $[t^h]N_p$ has $x$-degree less than $2p$ for those $h$.
Extracting $x^{n+2p}$ proves the asserted range $n\geq0$.
For $h=0$ the degree bound $2p$ gives $n\geq1$ instead.
\end{proof}

\begin{corollary}[Parity and the three-hit conjecture]\label{cor:parity}
Every odd-hit coefficient $u_{2j+1}(n)$ is even.
Theorem~\ref{thm:cong} holds, as does the two-hit congruence for every prime.
\end{corollary}
\begin{proof}
Modulo $2$, only $k=0,1$ survive in \eqref{eq:bivariate}, giving
\[
 \sum_{n\geq0}U_n(t)x^n
 =1+\frac{x}{1+x^2+x^2t^2},\qquad
 C_2(x)=\frac{x^3}{1+x^4}.
\]
The first identity gives odd-hit parity. The second gives
$(1+x^2)^2C_2=x^3$, which proves the two-hit congruence also at $p=2$.
For three hits, use parity at $p=2$ and Theorem~\ref{thm:uniform} at $p>3$.
\end{proof}

The extension to $p=3$ is not asserted in A000159 and fails at $n=4$:
$a(4)+2a(7)+a(10)=697756\equiv1\pmod3$.

\begin{corollary}[Prime-square antiperiods]\label{cor:antiperiod}
For every odd prime $p$ and $n\geq0$,
$a(n+p^2)\equiv-a(n)\pmod p$.
\end{corollary}
\begin{proof}
For $p>3$, set $v_k=a(n+kp)$ in $\mathbb F_p$.
Theorem~\ref{thm:cong} gives $v_{k+2}+2v_{k+1}+v_k=0$, so induction gives
\[
 v_k=(-1)^k\{v_0-k(v_0+v_1)\}.
\]
Take $k=p$. For $p=3$, extract $t^3$ from the three terms $k=0,1,2$
in \eqref{eq:bivariate}; direct expansion gives
\[
 A(x)=\frac{x^7-x^5-x^3}{(1+x)^8},\qquad
 (1+x^9)A=x^8+x^7-x^6-x^5-x^4-x^3
 \quad\text{in }\mathbb F_3[[x]].
\]
The polynomial on the right has degree $8$, proving the result for $p=3$.
\end{proof}

\begin{corollary}[Hit diagonals]\label{cor:diagonals}
The prime-shift conjectures in A000033 and A000181 hold on their full
stated domains. So does the corresponding assertion in A000185.
For A000179 and every prime $p$,
\[
 u_0(n+4p)-2u_0(n+2p)+u_0(n)\equiv0\pmod p\qquad(n\geq1).
\]
The same uniform prime-shift result holds for A058089 at $p>6$ and for
A058090 at $p>7$, on their source domains.
\end{corollary}
\begin{proof}
The diagonals A000033, A000181 and A000185 are $u_2,u_4,u_5$.
Use Theorem~\ref{thm:uniform} when $p>h$ and parity when $p=2$ and $h$ is odd;
Corollary~\ref{cor:parity} covers the remaining two-hit case $p=2$.
For zero hits, multiply the annihilator $(E^p+1)^2$ by $(E^p-1)^2$,
where $Ef(n)=f(n+1)$. Their product is $E^{4p}-2E^{2p}+1$.
Finally A058089 and A058090 are $u_6$ and $u_7$.
\end{proof}

\begin{corollary}[The two-hit quotient]\label{cor:quotient}
Let $q(n)=u_2(n)/n$ for $n\geq1$ and $q(0)=0$.
These are integers, equal to A000426, and for every prime $p$ and $n\geq1$,
\[
 q(n)+2q(n+p)+q(n+2p)\equiv0\pmod p.
\]
\end{corollary}
\begin{proof}
Define the integral formal series
\[
 Q(x)=\sum_{k\geq0}(2k+1)k!\frac{x^{k+2}}{(1+x)^{2k+2}}.
\]
Differentiating its $k$th term gives
$(2k+1)k!x^{k+2}\{k+2-kx\}/(1+x)^{2k+3}$ after multiplication by $x$.
This is exactly $[t^2]$ of the $k$th term of \eqref{eq:bivariate}.
Thus $xQ'=C_2$, proving both integrality and the quotient identity.
Modulo $p$, multiplication of the truncated series by $(1+x)^{2p}$
gives a polynomial of degree at most $2p$.
Coefficient extraction at $n+2p$ proves the congruence for $n\geq1$.
\end{proof}

\section{Classical identities and the corrected index}
Let
\[
 H_m(t)=\sum_{j=0}^m\binom{2m-j}{j}(m-j)!(t-1)^j,
 \qquad v(m)=H_m(0)=\mathrm{A000271}(m).
\]
Kaplansky--Riordan \cite[p.~122, equation (33)]{kr} give
\begin{equation}\label{eq:classical}
 U_n'=2nH_{n-1},\qquad
 U_n'=2nU_{n-1}-\frac{n(t-1)}{n-1}U_{n-1}'\quad(n\geq2).
\end{equation}
The first identity also holds at $n=1$ by direct calculation.

\begin{corollary}[Classical one-hit and two/three-hit identities]\label{cor:classical}
For $n\geq1$, $u_1(n)=2n v(n-1)$.
For $n\geq3$,
\[
 3\{n a(n+1)-(n+1)a(n)\}=2(n^2-1)u_2(n).
\]
These identities are consequences of Kaplansky--Riordan (1946).
\end{corollary}
\begin{proof}
Evaluate the first identity in \eqref{eq:classical} at $t=0$.
For the second, replace $n$ by $n+1$ in the other identity and extract $t^2$:
$3a(n+1)=2(n+1)u_2(n)-(n+1)\{2u_2(n)-3a(n)\}/n$.
Multiplication by $n$ gives the result.
\end{proof}

A000425's third FORMULA line (Mathar, 2 November 2015) instead conjectures
\[
u_1(n)=2n\mathrm{A000271}(n).
\]
At $n=4$ the two sides are $8$ and $24$.
The preceding corollary gives the corrected shift, with its classical proof.
The second identity is the A000033 relation recorded by van Hoeij on
9 June 2019, now explicitly attributed to the earlier identity.

\section{Exact recurrence certificates for sibling entries}\label{sec:siblings}
Here is a uniform algebraic way to check the additional source recurrences.
Write $\theta=z\,d/dz$ and $\theta^jF=f_jF+g_j$.
Starting from $(f_0,g_0)=(1,0)$, the factorial-series equation gives
\[
 f_{j+1}=zf_j'+(1/z-1)f_j,\qquad
 g_{j+1}=zg_j'-f_j/z.
\]
For $h\geq1$ put
$\alpha_h(k)=\binom{2k+h}{h}$,
$\beta_h(k)=\binom{2k+h-1}{h-1}$, and put
$\alpha_0=1,\beta_0=0$. Extracting $t^h$ from
\eqref{eq:bivariate} yields
\begin{equation}\label{eq:allmodels}
 C_h=\frac{x^h}{(1+x)^{h+1}}
 \left[\{(1-x)\alpha_h(\theta)+(1+x)\beta_h(\theta)\}F\right]_{z=Z}.
\end{equation}
The recursions for $f_j,g_j$ therefore give explicit rational functions
$(R_h,S_h)$ with $C_h=R_hB+S_h$.
Differentiation of a pair is the exact rational operation
\[
 \partial(R,S)=(R'+RT,S'+RV),\qquad
 T=\frac{1-x^3}{x^2(1+x)},\quad V=\frac{x^2-1}{x^2}.
\]
To check a source recurrence $\sum_{j=0}^m d_j(n)u_h(n-j)=0$, apply
\begin{equation}\label{eq:sourceoperator}
 \mathcal L=\sum_{j=0}^m x^jd_j(D+j),\qquad D=x\partial,
\end{equation}
to $(R_h,S_h)$. Substitution gives the following rational identities;
the complete coefficient arrays and the operations above are in
\texttt{models.py}.
\begin{center}\small
\begin{tabular}{llll}
\toprule
Source recurrence & $h$ & $\mathcal L(C_h)$ & Source-domain range\\
\midrule
A000425, FORMULA 2 &1&$0$&$n\geq4$\\
A000033, FORMULA 4 &2&$0$&$n\geq5$\\
A000159, Maplesoft reduction &3&$0$&$n\geq6$\\
A000181, FORMULA 1 &4&$0$&$n\geq8$\\
A000185, FORMULA 1 &5&$-3202597896x^5$&$n\geq11$\\
A000185, FORMULA 2 &5&$0$&$n\geq9$\\
\bottomrule
\end{tabular}\end{center}
Each identity has zero coefficient of $B$; the displayed remainder
is an exact polynomial. Equations \eqref{eq:allmodels} and
\eqref{eq:sourceoperator}, or polynomial expansion after clearing
denominators, verify the table without fitting any initial terms.
For $n\geq m$ the coefficient of $x^n$ is the required recurrence sum.
The remainder vanishes in all the listed ranges. This proves the
recurrences there, including the conjectures of Mathar in A000425,
A000181 and A000185.

The official Maplesoft example \cite{maple} starts from Mathar's
recurrence and explicitly says ``Assuming the input R is correct''.
Its third-order reduction is thus conditional on that input;
Theorem~\ref{thm:rec} and the identity above supply the coefficient-defined
sequence's recurrence premise. No earlier unconditional proof of this
coefficient bridge is inferred from the software example.

\section{Source coverage and verification}
The following locators identify the precise source statements
\cite{a159,a58087,siblings}. A FORMULA number counts the displayed
FORMULA lines in the indicated entry. Dates are contribution dates.
\begingroup\small
\begin{longtable}{p{.46\linewidth}p{.45\linewidth}}
\toprule Source item & Coverage\\\midrule\endhead
A000159 rev.~54, FORMULA 1, Mathar, 2 Nov.~2015
& Proved: recurrence for $n\geq9$; zero extension for $n=6,7,8$.\\[3pt]
A000159 rev.~54, FORMULA 3, van Hoeij, 10 June 2019
& Proved: prime-shift congruence, extended to $n\geq0$.\\[3pt]
A000159 FORMULA 2; A000425 FORMULA 1; A000033 and A000426 Bessel formulas
& Not covered: the analytic rounding error bounds.\\[3pt]
A000425 rev.~34, FORMULA 2, Mathar, 2 Nov.~2015
& Proved: recurrence for $n\geq4$ (Section~\ref{sec:siblings}).\\[3pt]
A000425 rev.~34, FORMULA 3, Mathar, 2 Nov.~2015
& Refuted-and-corrected: $8\ne24$ at $n=4$; corrected shift reproved; prior proof Kaplansky--Riordan, 1946, (33).\\[3pt]
A000033 rev.~60, FORMULA 4, Kotesovec, 26 Oct.~2012
& Proved: the stated recurrence, $n\geq5$.\\[3pt]
A000033 rev.~60, two/three-hit identity, van Hoeij, 9 June 2019
& Reproved; prior proof Kaplansky--Riordan, 1946, (33).\\[3pt]
A000033 rev.~60, last FORMULA, van Hoeij, 10 June 2019
& Proved: congruence for all primes, $n\geq1$; also at $n=0$ by zero extension.\\[3pt]
A000181 rev.~30, FORMULAS 1--2, Mathar, 2 Nov.~2015; van Hoeij, 10 June 2019
& Proved: recurrence for $n\geq8$ and congruence for $p>3$ on $n\geq4$.\\[3pt]
A000185 rev.~28, FORMULAS 1--3, Mathar, 2 Nov.~2015; van Hoeij, 13 June 2019
& Proved: recurrences for $n\geq11$ and $n\geq9$; congruence for $p\ne3,5$ on $n\geq5$.\\[3pt]
A000426 rev.~44, last FORMULA, van Hoeij, 13 June 2019
& Proved: quotient congruence for all primes, $n\geq1$.\\[3pt]
A000179 rev.~274, van Hoeij, 11 Jan.~2022, four-prime-shift assertion
& Proved for every prime, $n\geq1$, with the source's $u_0(1)=-1$.\\[3pt]
A058089 rev.~11 and A058090 rev.~14, defining diagonals
& Not posed; proved: uniform congruence for $p>6$ and $p>7$, respectively.\\[3pt]
A000159 with $p=3$
& Not asserted in the source; fails for $n=4$; prime-square antiperiod instead proved.\\[3pt]
A058087 rev.~66, FORMULAS 1--2; A094314 rev.~35
& Prior generating function (Jovovi\'c, 2009) and hit polynomial (Kaplansky--Riordan, 1946) used. A094314 agrees for $n\geq2$; its $n=1$ convention differs.\\[3pt]
Other linked asymptotic, nonlinear or negative-index formulas
& Not covered.\\
\bottomrule
\end{longtable}\endgroup

The Lean declaration
\href{https://github.com/the-omega-institute/trureturing/blob/8acba483bc80695fef00c6b8c491e2ae16d5145f/D5/S1/Recurrence/Invariants/MatharA000159Recurrence.lean#L656}{\texttt{MatharA000159Recurrence.result}}
states the recurrence for the rational coefficient
$a(n)=\coeff{t}{3}U_n(t)$ for all $n\geq9$.
Its definitions \texttt{U}, \texttt{a}, and \texttt{claim} are in the same
pinned file. This declaration covers Theorem~\ref{thm:rec} for $n\geq9$,
not the modular or sibling results. The zero extension, uniform congruences,
antiperiods, quotient theorem and classical deductions have the written
proofs above. The recurrence certificates are exact rational identities.
The accompanying \texttt{verify.py} checks these rational identities and
independently compares the defining rook sum with the recurrences and
congruences on the finite ranges specified in \texttt{README.md}.
These finite comparisons are separate from the all-index proofs.

\begin{thebibliography}{9}
\bibitem{a159}
The OEIS Foundation Inc., \emph{The On-Line Encyclopedia of Integer
Sequences}, A000159, revision 54, 21 June 2026.
\url{https://oeis.org/A000159/internal}.
FORMULA contributions by R. J. Mathar (2 November 2015) and Mark
van Hoeij (9 and 10 June 2019).
\bibitem{a58087}
The OEIS Foundation Inc., \emph{The On-Line Encyclopedia of Integer
Sequences}, A058087, revision 66, 5 November 2025.
\url{https://oeis.org/A058087/internal}.
Generating function recorded by Vladeta Jovovi\'c, 14 December 2009;
M\'enage hit polynomial formula recorded by William P. Orrick,
9 August 2020.
\bibitem{kr}
I. Kaplansky and J. Riordan,
\emph{The probl\`eme des m\'enages},
Scripta Mathematica \textbf{12} (1946), 113--124,
especially pp.~121--122, equations (25), (26) and (33).
\url{https://oeis.org/A000166/a000166_1.pdf}.
\bibitem{siblings}
The OEIS Foundation Inc., entries
\href{https://oeis.org/A000179/internal}{A000179} (rev.~274),
\href{https://oeis.org/A000425/internal}{A000425} (rev.~34),
\href{https://oeis.org/A000033/internal}{A000033} (rev.~60),
\href{https://oeis.org/A000271/internal}{A000271} (rev.~127),
\href{https://oeis.org/A000426/internal}{A000426} (rev.~44),
\href{https://oeis.org/A000181/internal}{A000181} (rev.~30),
\href{https://oeis.org/A000185/internal}{A000185} (rev.~28),
\href{https://oeis.org/A058089/internal}{A058089} (rev.~11),
\href{https://oeis.org/A058090/internal}{A058090} (rev.~14), and
\href{https://oeis.org/A094314/internal}{A094314} (rev.~35).
Contribution locators are specified in the coverage table.
\bibitem{maple}
Maplesoft, \emph{LREtools[MinimalRecurrence]}, Examples, outputs (3)--(7)
and the paragraph beginning ``Assuming the input R is correct'';
command introduced in Maple 2021.
\url{https://www.maplesoft.com/support/help/Maple/view.aspx?path=LREtools%2FMinimalRecurrence}.
\end{thebibliography}
\bigskip
\noindent\footnotesize
\copyright\ 2026 Haobo Ma and Wenlin Zhang.
Licensed under Creative Commons Attribution 4.0 International
(CC BY 4.0).
\end{document}
